%!TEX program = xelatex
\documentclass[aspectratio=169,fontset=none]{ctexbeamer}
\setCJKmainfont[
  Path=/System/Library/Fonts/,
  UprightFont={STHeiti Light.ttc},
  BoldFont={STHeiti Medium.ttc}
]{STHeiti Light.ttc}
\setCJKsansfont[
  Path=/System/Library/Fonts/,
  UprightFont={STHeiti Light.ttc},
  BoldFont={STHeiti Medium.ttc}
]{STHeiti Light.ttc}
\setCJKmonofont[
  Path=/System/Library/Fonts/Supplemental/
]{Songti.ttc}

% 引用科大专属主题包
\usepackage[bluetheme]{ustcbeamer}
\input{ustctheme.tex}
\usepackage{graphicx}
\usepackage{booktabs}
\usepackage{hyperref}
\usepackage{tikz}
\usetikzlibrary{positioning,fit,arrows.meta,shapes.geometric}

% ---------------------------------------------------
% 封面信息配置
% ---------------------------------------------------
\title[无线腱反射采集系统]{
    无线腱反射诱发与多传感器同步采集系统
}
\subtitle{组会汇报}
\author[胡彦卓]{报告人：胡彦卓}
\institute[USTC]{
    微电子学院，神经肌肉控制实验室
}
\date{2026年5月7日}

\begin{document}

% ---------------------------------------------------
\maketitleframe

% ---------------------------------------------------
\begin{frame}{汇报大纲}
    \tableofcontents
\end{frame}

% ---------------------------------------------------
\AtBeginSection[]{
\setbeamertemplate{footline}[footlineoff]
  \begin{frame}
    \frametitle{大纲}
    \tableofcontents[currentsection,subsectionstyle=show/show/hide]
  \end{frame}
\setbeamertemplate{footline}[footlineon]
}

% ===================================================
\section{系统目标与总体架构}

\begin{frame}{系统目标：同步记录腱反射过程中的关键事件}
  \begin{columns}
    \begin{column}{0.52\textwidth}
      \begin{block}{实验对象}
        \begin{itemize}
          \item 膝跳反射：叩击髌腱后，小腿产生快速摆动。
          \item 目标是把叩击、肌电响应和小腿运动放到统一时间轴上。
        \end{itemize}
      \end{block}
      \begin{block}{需要同步采集的信号}
        \begin{itemize}
          \item 肌电：反映反射诱发后的神经肌肉响应。
          \item 小腿角速度：刻画膝跳反射时小腿摆动过程。
          \item 锤体加速度：捕捉叩击时刻与叩击波形。
        \end{itemize}
      \end{block}
    \end{column}
    \begin{column}{0.48\textwidth}
      \centering
      \begin{tikzpicture}[
        scale=0.76,
        transform shape,
        node distance=0.36cm,
        sensor/.style={draw, rounded corners=2pt, align=center, minimum width=3.5cm, minimum height=0.72cm, fill=white},
        event/.style={draw, rounded corners=2pt, align=center, minimum width=3.9cm, minimum height=0.82cm, fill=blue!5},
        arrow/.style={-{Stealth[length=2.2mm]}, thick}
      ]
        \node[event] (tap) {叩击髌腱\\锤体加速度};
        \node[event, below=of tap] (emg) {诱发肌电\\ADS1298};
        \node[event, below=of emg] (motion) {小腿摆动\\角速度};
        \node[sensor, right=0.75cm of emg] (axis) {统一时间轴\\反射潜伏期与运动响应};
        \draw[arrow] (tap.east) -- (axis.west);
        \draw[arrow] (emg.east) -- (axis.west);
        \draw[arrow] (motion.east) -- (axis.west);
      \end{tikzpicture}
      \vspace{0.22cm}

      \begin{block}{核心要求}
        多节点、无线化、时间同步、可联调。
      \end{block}
    \end{column}
  \end{columns}
\end{frame}

\begin{frame}{总体架构：一主两从的同步采集链路}
  \centering
  \vspace{0.08cm}
  \begin{tikzpicture}[
    scale=0.88,
    transform shape,
    box/.style={draw, rounded corners=2pt, align=center, minimum width=3.0cm, minimum height=0.86cm, fill=white},
    master/.style={draw, rounded corners=2pt, align=center, minimum width=3.45cm, minimum height=1.05cm, fill=blue!7},
    pc/.style={draw, rounded corners=2pt, align=center, minimum width=3.15cm, minimum height=0.96cm, fill=green!7},
    espnow/.style={-{Stealth[length=2.2mm]}, thick, teal!75!black},
    beacon/.style={-{Stealth[length=2.2mm]}, thick, dashed, orange!85!black},
    usb/.style={-{Stealth[length=2.2mm]}, thick, blue!75!black},
    tag/.style={font=\scriptsize, fill=white, inner sep=1.2pt}
  ]
    \node[box] (slave1) at (-3.5,1.10) {从机 1\\MPU6500 角速度\\小腿摆动};
    \node[box] (slave2) at (-3.5,-1.10) {从机 2\\MPU6500 加速度\\叩击锤};
    \node[master] (master) at (0.85,0) {主机 Master\\ADS1298 肌电采集\\时间基准与数据汇聚};
    \node[pc] (pc) at (5.0,0) {上位机\\实时显示\\CSV 保存};

    \draw[espnow] (slave1.east) -- (master.west);
    \draw[espnow] (slave2.east) -- (master.west);
    \draw[beacon, shorten >=2pt] (master.165) to[out=165, in=0] (slave1.0);
    \draw[beacon, shorten >=2pt] (master.195) to[out=195, in=0] (slave2.0);
    \draw[usb] (master.east) -- (pc.west);

    \fill[teal!75!black] (-3.38,-2.48) rectangle (-3.20,-2.36);
    \node[tag, anchor=west] at (-3.10,-2.42) {ESP-NOW 批量包};
    \fill[orange!85!black] (0.08,-2.48) rectangle (0.26,-2.36);
    \node[tag, anchor=west] at (0.36,-2.42) {Beacon：Master $\rightarrow$ Slave};
    \fill[blue!75!black] (4.08,-2.48) rectangle (4.26,-2.36);
    \node[tag, anchor=west] at (4.36,-2.42) {USB CDC};
  \end{tikzpicture}

  \vspace{0.12cm}
  \begin{block}{数据组织方式}
    \begin{itemize}
      \item Master 以自身时间为全局时间基准，Slave 收到 Beacon 后将本地采样时间换算到 Master 时间轴。
      \item Master 将本地肌电、从机 IMU 与诊断信息统一封装后发给上位机。
    \end{itemize}
  \end{block}
\end{frame}

\begin{frame}{节点分工}
  \vspace{0.08cm}
  \begin{table}
    \centering
    \scriptsize
    \begin{tabular}{p{0.13\textwidth}p{0.18\textwidth}p{0.13\textwidth}p{0.42\textwidth}}
      \toprule
      节点 & 传感器/执行器 & 安装位置 & 主要作用 \\
      \midrule
      Master & ADS1298 & 肌电电极侧 & 采集诱发肌电，广播时间基准，汇聚数据 \\
      Slave-01 & MPU6500 陀螺仪 & 小腿 & 采集膝跳反射时小腿角速度 \\
      Slave-02 & MPU6500 加速度计 & 叩击锤 & 捕捉叩击时刻和叩击波形 \\
      Motor & 电机 + 驱动模块 & 叩击执行机构 & 后续通过蓝牙控制叩击动作 \\
      \bottomrule
    \end{tabular}
  \end{table}

  \vspace{0.22cm}
  \begin{columns}
    \begin{column}{0.50\textwidth}
      \begin{block}{采集链路}
        Master 与两个 Slave 构成同步采集系统，数据最终由 Master 输出到上位机。
      \end{block}
    \end{column}
    \begin{column}{0.50\textwidth}
      \begin{block}{控制链路}
        Motor 节点作为独立蓝牙控制对象，后续与采集链路做触发和时序联调。
      \end{block}
    \end{column}
  \end{columns}
\end{frame}

% ===================================================
\section{同步采集方案}

\begin{frame}{时间同步与数据流}
  \begin{columns}
    \begin{column}{0.48\textwidth}
      \begin{block}{同步思路}
        \begin{itemize}
          \item Master 周期广播带时间戳的 Beacon。
          \item Slave 根据 Beacon 估算主从时间偏移。
          \item Slave 每个采样点都写入换算后的 Master 时间戳。
          \item 上位机按统一时间戳对齐肌电、角速度和加速度。
        \end{itemize}
      \end{block}
    \end{column}
    \begin{column}{0.52\textwidth}
      \centering
      \begin{tikzpicture}[
        x=0.92cm,y=0.92cm,
        scale=0.92,
        transform shape,
        line/.style={thick},
        arrow/.style={-{Stealth[length=2.1mm]}, thick},
        label/.style={font=\scriptsize},
        tag/.style={font=\tiny, fill=white, inner sep=1pt}
      ]
        \draw[line] (0,2.4) -- (5.4,2.4);
        \draw[line] (0,1.4) -- (5.4,1.4);
        \draw[line] (0,0.4) -- (5.4,0.4);
        \node[label, left] at (0,2.4) {Master};
        \node[label, left] at (0,1.4) {Slave};
        \node[label, left] at (0,0.4) {PC};

        \foreach \x in {0.6,1.6,2.6,3.6,4.6}{
          \draw[blue, thick] (\x,2.28) -- (\x,2.52);
        }
        \foreach \x in {0.9,1.9,2.9,3.9,4.9}{
          \draw[teal!80!black, thick] (\x,1.28) -- (\x,1.52);
        }
        \draw[arrow, dashed] (1.05,2.28) -- (1.05,1.52);
        \draw[arrow, dashed] (3.55,2.28) -- (3.55,1.52);
        \draw[arrow] (2.05,1.52) .. controls (2.28,1.95) and (2.62,2.02) .. (2.95,2.28);
        \draw[arrow] (4.05,1.52) .. controls (4.28,1.92) and (4.60,1.98) .. (4.95,2.28);
        \draw[arrow] (5.25,2.28) -- (5.25,0.52);
        \node[tag, anchor=west] at (0.55,-0.18) {虚线：Beacon};
        \node[tag, anchor=west] at (2.05,-0.18) {弧线：IMU batch};
        \node[tag, anchor=west] at (3.95,-0.18) {竖线：CDC batch};
      \end{tikzpicture}
    \end{column}
  \end{columns}

  \vspace{0.16cm}
  \begin{block}{关键点}
    采样、无线通信和串口输出分任务运行，并用 CRC 对 ESP-NOW 与 USB CDC 数据帧做校验。
  \end{block}
\end{frame}

\begin{frame}{上位机视角：多源信号进入同一数据表}
  \begin{columns}
    \begin{column}{0.48\textwidth}
      \begin{block}{统一输出}
        \begin{itemize}
          \item Master 将各节点样本打包为 USB CDC 批量帧。
          \item 上位机解析帧头、类型、样本数、节点来源和 CRC。
          \item 数据保存为 CSV，便于后续延迟、波形和幅值分析。
        \end{itemize}
      \end{block}
    \end{column}
    \begin{column}{0.52\textwidth}
      \begin{table}
        \centering
        \small
        \scriptsize
        \begin{tabular}{p{0.18\linewidth}p{0.24\linewidth}p{0.38\linewidth}}
          \toprule
          source & 信号 & 用途 \\
          \midrule
          Master & EMG & 判断反射诱发肌电响应 \\
          Slave-01 & Gyro $x/y/z$ & 分析小腿摆动起始和速度变化 \\
          Slave-02 & Accel $x/y/z$ & 定位叩击瞬间和叩击强度 \\
          Diag & 状态信息 & 检查丢包、CRC、同步状态 \\
          \bottomrule
        \end{tabular}
      \end{table}
      \vspace{0.1cm}
      \begin{block}{分析目标}
        从锤体加速度峰值出发，比较肌电响应潜伏期与小腿运动延迟。
      \end{block}
    \end{column}
  \end{columns}
\end{frame}

% ===================================================
\section{当前进展}

\begin{frame}{当前进展：主从机同步采集已完成}
  \begin{columns}
    \begin{column}{0.50\textwidth}
      \begin{block}{已经完成}
        \begin{itemize}
          \item 完成 Master、Slave-01、Slave-02 三类采集固件。
          \item 完成 Master 对两个 Slave 的时间同步和数据汇聚。
          \item 完成 ESP-NOW 批量传输与 USB CDC 批量输出。
          \item 完成上位机实时接收、显示和 CSV 保存流程。
        \end{itemize}
      \end{block}
    \end{column}
    \begin{column}{0.50\textwidth}
      \begin{block}{目前可以验证}
        \begin{itemize}
          \item 肌电、角速度、加速度可以进入同一采集流程。
          \item Slave 数据带有换算后的 Master 时间戳。
          \item 上位机可用于观察同步波形和导出实验数据。
        \end{itemize}
      \end{block}
      \begin{block}{阶段结论}
        采集链路已经具备后续整机联调的基础。
      \end{block}
    \end{column}
  \end{columns}
\end{frame}

\begin{frame}{当前系统状态}
  \centering
  \begin{tikzpicture}[
    stage/.style={draw, rounded corners=2pt, align=center, minimum width=3.05cm, minimum height=0.88cm, fill=white},
    done/.style={draw, rounded corners=2pt, align=center, minimum width=3.05cm, minimum height=0.88cm, fill=green!10},
    todo/.style={draw, rounded corners=2pt, align=center, minimum width=3.05cm, minimum height=0.88cm, fill=orange!10},
    arrow/.style={-{Stealth[length=2.2mm]}, thick}
  ]
    \node[done] (emg) {主机肌电采集\\已接入};
    \node[done, right=0.42cm of emg] (imu) {双从机 IMU\\已同步采集};
    \node[done, right=0.42cm of imu] (host) {上位机接收\\已打通};
    \node[todo, right=0.42cm of host] (motor) {蓝牙叩击控制\\待联调};
    \draw[arrow] (emg) -- (imu);
    \draw[arrow] (imu) -- (host);
    \draw[arrow] (host) -- (motor);
  \end{tikzpicture}

  \vspace{0.35cm}
  \begin{columns}
    \begin{column}{0.50\textwidth}
      \begin{block}{采集侧}
        重点问题已经从“能否同步采到多源数据”转向“整机实验时序是否稳定”。
      \end{block}
    \end{column}
    \begin{column}{0.50\textwidth}
      \begin{block}{执行侧}
        蓝牙控制叩击将作为下一阶段引入，减少手动叩击的不确定性。
      \end{block}
    \end{column}
  \end{columns}
\end{frame}

% ===================================================
\section{下一步计划}

\begin{frame}{下一步计划：蓝牙控制叩击}
  \begin{columns}
    \begin{column}{0.52\textwidth}
      \begin{block}{目标}
        \begin{itemize}
          \item 使用蓝牙控制叩击执行节点，替代人工手动叩击。
          \item 保持叩击动作和采集链路相互独立，降低系统耦合。
          \item 通过锤体加速度反推真实叩击时刻，保证分析以实测事件为准。
        \end{itemize}
      \end{block}
    \end{column}
    \begin{column}{0.48\textwidth}
      \centering
      \begin{tikzpicture}[
        node distance=0.42cm,
        box/.style={draw, rounded corners=2pt, align=center, minimum width=3.7cm, minimum height=0.82cm, fill=white},
        arrow/.style={-{Stealth[length=2.2mm]}, thick}
      ]
        \node[box] (pc) {上位机\\蓝牙控制指令};
        \node[box, below=of pc] (motor) {Motor 节点\\PWM 控制};
        \node[box, below=of motor] (hammer) {叩击机构\\产生可控叩击};
        \node[box, below=of hammer] (acc) {锤体加速度\\记录实际叩击波形};
        \draw[arrow] (pc) -- node[right, font=\scriptsize]{BLE} (motor);
        \draw[arrow] (motor) -- (hammer);
        \draw[arrow] (hammer) -- (acc);
      \end{tikzpicture}
    \end{column}
  \end{columns}
\end{frame}

\begin{frame}{下一步计划：整机联调测试}
  \begin{columns}
    \begin{column}{0.50\textwidth}
      \begin{block}{联调内容}
        \begin{itemize}
          \item 蓝牙叩击控制与采集上位机流程联调。
          \item 验证叩击加速度、肌电响应、小腿角速度之间的时序关系。
          \item 观察连续测试下的丢包、同步漂移和波形稳定性。
        \end{itemize}
      \end{block}
    \end{column}
    \begin{column}{0.50\textwidth}
      \begin{block}{预期输出}
        \begin{itemize}
          \item 得到包含叩击、肌电、运动响应的完整实验数据。
          \item 形成初步的反射潜伏期与运动响应分析流程。
          \item 根据联调结果继续优化固定方式、采样参数和控制节拍。
        \end{itemize}
      \end{block}
    \end{column}
  \end{columns}

  \vspace{0.22cm}
  \begin{block}{近期里程碑}
    完成“蓝牙触发叩击 -- 加速度捕捉叩击 -- 肌电与角速度同步记录”的闭环验证。
  \end{block}
\end{frame}

\begin{frame}{}
  \vspace{0.30\textheight}
  \centering
  \Huge 谢谢！
\end{frame}

\end{document}
